\section{Consequences, checks, and scope}

The main theorem is a realization result.  It does more than produce a large
collection of examples: given a finite \'etale algebra, it builds a Keller map
whose full fiber is that algebra.  Several consequences are then immediate.
A squarefree polynomial of degree at least three becomes an exact inverse
polynomial after one translation; an irreducible polynomial gives a connected
fiber; and a degree-$N$ input is realized in $\A^3$ with coordinate degree at
most $6N+2$.

More importantly, the construction does not forget the arithmetic structure.
If the original algebra has a particular real signature, splitting field,
Galois action, or collection of local factorizations, the fiber has the same
data.  The same applies to local and rational points.  Thus an intersective
polynomial with no rational root gives a Hasse-failing Keller fiber directly,
without adding auxiliary factors or extra variables.

This is a transfer statement, not an inverse Galois theorem.  It does not claim
that arbitrary local specifications can always be globalized.  It says that
once a finite \'etale algebra with the desired arithmetic properties exists,
the Keller construction preserves it exactly.

The realization is also far from unique.  Other constructions in the
accompanying repository place the same finite \'etale algebra inside weighted,
cancellation, and quadratic-gauge Keller maps that need not be stably
equivalent.  That is a different question, and I leave it for a separate
paper.

\paragraph{Relation to earlier work.}
To my knowledge, the exact rank classification proved here has not appeared
before.  I searched for earlier results up to 24 July 2026, but this is
necessarily a qualified statement: searches can miss unpublished work or a
theorem stated in different language.

The closest constructions have a different scope.  Gallagher gives explicit
Keller maps in every geometric degree together with complete split rational
fibers \cite{Gallagher}; these realize $K^N$, rather than an arbitrary finite
\'etale $K$-algebra.  Miranda--Neto studies radicality and maximal-ideal
decompositions of fibers that already come from Keller maps \cite{Miranda}.
Lipton and Markakis connect rational images with Hilbert irreducibility
\cite{LiptonMarkakis}, and Shaska studies graded Keller maps and arithmetic
thinness \cite{Shaska}.  None of these results states the prescribed-fiber
realization theorem used here.

\paragraph{Independent computer checks.}
The formulas are uniform in the degree, but expanding them by hand is not a
pleasant way to find a sign error.  I therefore checked the construction in
three independent ways.

The SymPy script
\[
 \texttt{scripts/verify\_universal\_quadratic\_gauge.py}
\]
verifies the source-chart identity, the marked-line formulas, the Jacobian
cancellation, every generic coefficient contribution, and the degree bound.
These are symbolic identities with arbitrary coefficients, not numerical
tests at a few sample points.

The Singular script
\[
 \texttt{scripts/verify\_universal\_quadratic\_gauge.sing}
\]
starts again over a rational function field in six independent coefficients.
It expands the complete degree-six map, obtains determinant exactly $-2$, and
checks the degree profile $(7,38,36)$.  This gives a separate CAS
implementation of the same construction.

Finally,
\[
 \texttt{scripts/verify\_finite\_etale\_keller\_fibers.py}
\]
builds translated seeds in degrees three, four, and five, expands the actual
maps, reconstructs the prescribed quotient rings, and checks the explicit
rank-five Hasse example and the fixed-map family.  These bounded examples are
regression tests.  The proof for arbitrary degree is the coefficientwise
argument in Proposition~\ref{prop:gauge} and
Subsection~\ref{subsec:gauge-identities}.

\paragraph{Lean formalization.}
The Lean project lives in
\[
 \texttt{formal/finite-etale-keller}.
\]
It contains no \texttt{sorry} and introduces no project-specific axioms.  The
formal development covers the constructive core of the paper.  Given a
squarefree polynomial $P$ of degree at least three, it:
\begin{enumerate}
\item proves that an admissible translation exists and chooses one;
\item constructs the complete all-degree map as an actual
      $\operatorname{MvPolynomial}(\operatorname{Fin}3,K)$ expression;
\item proves that its normalized Jacobian is one and that all coordinate
      degrees are at most $6\deg P+2$;
\item proves that the inverse polynomial at the selected target is
      $P(a+S)$; and
\item identifies the literal map fiber naturally with
      $\operatorname{Hom}_{K\text{-alg}}(K[T]/(P),A)$ for every commutative
      test algebra $A$.
\end{enumerate}

The final construction and degree certificate are exposed as
\texttt{automaticRealizationMap\_certificate}; the natural fiber equivalence
is \texttt{automaticJacobianOneFiberRepresentingEquiv\_natural}.  Their axiom
reports contain only Lean's standard logical principles
\texttt{propext}, \texttt{Classical.choice}, and \texttt{Quot.sound}.

Lean does not currently cover every theorem used in the final classification.
The generic irreducibility argument for the function-field degree,
monogenicity of arbitrary finite \'etale products, and the
Campbell--Razar--Wright degree-two theorem remain ordinary mathematical proofs.
The Chebotarev, Dirichlet, Hensel, and prime-counting arguments in the
arithmetic applications are likewise outside the formal development.  I list
this boundary explicitly because a green Lean build should not be read as a
claim that those external inputs have also been formalized.

\paragraph{Provenance and AI assistance.}
The recent three-dimensional counterexample was publicly attributed to
Alp\"oge and Fable.  Gallagher independently developed a weighted all-degree
viewpoint and a public atlas of complete rational fibers \cite{Gallagher}.
The present proof uses a different, root-engineered quadratic gauge.  Berend
and Bilu supplied the classical intersective quintic, while Campbell, Razar,
and Wright supplied the degree-two obstruction
\cite{Campbell,Razar,Wright}.

I used generative AI tools during exploration, symbolic-checker development,
proof review, literature search, Lean formalization, and manuscript editing.
They were useful as fast but fallible collaborators.  I checked the stated
mathematical claims, rewrote the exposition, and take responsibility for the
paper.

\begin{thebibliography}{11}
\small

\bibitem{BerendBilu}
D.~Berend and Y.~Bilu,
\emph{Polynomials with roots modulo every integer},
Proc. Amer. Math. Soc. \textbf{124} (1996), 1663--1671,
\href{https://doi.org/10.1090/S0002-9939-96-03210-8}
{doi:10.1090/S0002-9939-96-03210-8}.

\bibitem{Campbell}
L.~A. Campbell,
\emph{A condition for a polynomial map to be invertible},
Math. Ann. \textbf{205} (1973), 243--248,
\href{https://doi.org/10.1007/BF01349234}
{doi:10.1007/BF01349234}.

\bibitem{Gallagher}
A.~Gallagher,
\emph{Counterexample atlas: generic fiber degrees 3 through 100},
online exact certificate atlas, 2026,
\href{https://jacobianfun.org/counterexamples}{jacobianfun.org}.

\bibitem{LiptonMarkakis}
R.~J. Lipton and E.~Markakis,
\emph{Some remarks on the Jacobian conjecture and connections with
Hilbert's irreducibility theorem},
\href{https://arxiv.org/abs/math/0507525}{arXiv:math/0507525}, 2005.

\bibitem{Miranda}
C.~B. Miranda--Neto,
\emph{An ideal-theoretic approach to Keller maps},
Proc. Edinburgh Math. Soc. \textbf{62} (2019), 1033--1044,
\href{https://doi.org/10.1017/S0013091519000099}
{doi:10.1017/S0013091519000099}.

\bibitem{Razar}
M.~Razar,
\emph{Polynomial maps with constant Jacobian},
Israel J. Math. \textbf{32} (1979), 97--106,
\href{https://doi.org/10.1007/BF02764906}
{doi:10.1007/BF02764906}.

\bibitem{Shaska}
T.~Shaska,
\emph{Graded Keller maps and the Jacobian conjecture},
\href{https://arxiv.org/abs/2607.20210}{arXiv:2607.20210}, 2026.

\bibitem{Wright}
D.~Wright,
\emph{On the Jacobian conjecture},
Illinois J. Math. \textbf{25} (1981), 423--440,
\href{https://doi.org/10.1215/ijm/1256047158}
{doi:10.1215/ijm/1256047158}.

\end{thebibliography}
